%% ===================================================================
%%  RF Signal Security Taxonomy -- tool paper
%%  SKELETON / GUIDE DRAFT. Not submission copy.
%%
%%  Every \guide{...} block is scaffolding addressed to the author:
%%  what belongs in that slot, what a reviewer will attack, what must
%%  be measured rather than asserted. Set \showguidesfalse to hide
%%  them all and read the prose on its own.
%%
%%  Build:  latexmk -pdf main.tex
%% ===================================================================

\documentclass[journal]{IEEEtran}

\usepackage{amsmath}
\usepackage{graphicx}
\usepackage{booktabs}
\usepackage{xcolor}
\usepackage{url}
\usepackage[hidelinks]{hyperref}

% ---- author-facing guidance ----------------------------------------
\definecolor{guidecol}{HTML}{8A5A00}
\newif\ifshowguides
\showguidestrue                 % <-- \showguidesfalse to hide all guide notes
\newcommand{\guide}[1]{%
  \ifshowguides
    {\par\smallskip\footnotesize\color{guidecol}%
     \textsf{\textbf{[guide]}~#1}\par\smallskip}%
  \fi}
\newcommand{\needcite}{\textcolor{red}{[CITE]}}
\newcommand{\needdata}{\textcolor{red}{[DATA]}}

\begin{document}

\title{SigVault-RF: An Interactive Relational Taxonomy \\ of RF Signal Security in Space Communications}

\author{%
\IEEEauthorblockN{Nicholas D. Redmond, Dipankar Dasgupta}\\
\IEEEauthorblockA{University of Memphis, Department of Computer Science\\
Memphis, TN, USA\\
Email: \{ndrdmond@memphis.edu, ddasgupt@memphis.edu\} }
}

\maketitle

% ====================================================================
\begin{abstract}
The literature on radio-frequency (RF) security is fragmented across terrestrial and orbital domains and often confined to the specific time it was written, moving slower than the attack surface is expanding. With AI pushing the frontiers of signal exploitation and serious hardware becoming cheaper and more accessible, the need for a modern and live taxonomy of RF Signal Security is obvious. We present SigVault, an interactive relational taxonomy in which knowledge of RF security becomes more accessible and engaging than ever before. Featured is a matrix of eleven signal-security facets as they relate to one another, rendered on a rotatble surface intended to map the literature in a way that is both navigatable and visually pleasing. Additionally, three coupled drill-down views are provided to allow users to visualize and interact with the details of the taxonomy in a memorable way. These views are named for their domain and purpose: TRANSEC, COMSEC, and our proprietary Signal Lab. This taxonomy is designed to reference the literature as well as map the research gaps in the field, shortening the time it takes to even find open problems. It is intended to be a living document that will continue to grow with the field which is essentialy in a time where AI advances are so rapid that the domain has changed before the academic journals can review and publish the findings.   
\end{abstract}

\begin{IEEEkeywords}·
Taxonomy, security visualization, radio-frequency security, satellite
communications, physical-layer security, matrix visualization, transmission
security.
\end{IEEEkeywords}

\guide{Trim the abstract to roughly 150--200 words before submission; the
version above is long on purpose so you can cut rather than invent. The
sentence that must survive intact is the one naming the contribution:
\emph{crossing} facets that prior work treats separately.}

% ====================================================================
\section{Introduction}
\IEEEPARstart{T}{he} security of a radio link is decided by properties
that no single research community owns. Whether a transmission can be
detected depends on its waveform and its antenna pattern; whether it can
be exploited depends on its framing and its encryption; whether it can be
denied depends on the geometry between transmitter, receiver, and
adversary. These properties are studied in electronic warfare, in
satellite systems engineering, in physical-layer security, and in mission
cyber defence, and each community has its own vocabulary, its own
canonical threats, and its own reference documents.



% ====================================================================
\section{Related Work}

\subsection{Structured Adversary Knowledge Bases}
SPARTA


\subsection{Surveys and Taxonomies of Satellite and PHY Security}


Two limitations recur. First, these treatments are largely
orbit-agnostic, which understates how strongly the threat model depends
on regime. Second, they are organised as prose surveys, so the knowledge
they contain cannot be queried along an axis the author did not
anticipate.

\subsection{Matrix and Adjacency Visualization}


% ====================================================================
\section{Taxonomy Structure}

\subsection{Facets}

\begin{table}[t]
\caption{The eleven facets and their families. Family membership drives
hue in the visual encoding.}
\label{tab:facets}
\centering
\footnotesize
\begin{tabular}{@{}lll@{}}
\toprule
\textbf{Facet} & \textbf{Family} & \textbf{Example members} \\
\midrule
Protocol     & System     & DVB-S2X, Starlink Ku, GPS L1 \\
Use-case     & System     & broadcast, TT\&C, backhaul \\
Band         & Spectrum   & HF, L, S, X, Ku/Ka, mmWave, optical \\
Modulation   & Spectrum   & BPSK, OQPSK, BOC(10,5), DSSS \\
Antenna      & Physical   & helical, parabolic, phased array \\
Propagation  & Physical   & vacuum, troposphere, multipath \\
Noise        & Physical   & thermal, co-channel, intentional \\
Attack       & Adversary  & intercept, follower jamming, meaconing \\
Objective    & Defence    & LPD, LPI, LPE, AJ, TFS \\
Mitigation   & Defence    & spreading, nulling, OSNMA \\
Limitation   & Residual   & PAPR cost, acquisition delay \\
\bottomrule
\end{tabular}
\end{table}


\subsection{Relations as the Unit of Knowledge}


\subsection{Evidence Density}


% ====================================================================
\section{Visual Design}
\label{sec:vis}

\subsection{Requirements}


\subsection{A Cyclic Axis}


\subsection{Encoding}


\subsubsection*{Rejected: identity as a colour blend}


\subsection{Interaction}




% ====================================================================
\section{System Architecture}

\subsection{One Graph, Four Projections}


\subsection{Provenance and Reproducibility}

\subsection{Implementation}

% ====================================================================
\section{Usage Scenarios}

\guide{Two scenarios is the right number for five pages; three if you cut
elsewhere. Each should be a task a real analyst has, not a feature tour,
and each should end with an answer the reader could not have obtained
from the sources cited in Section~II.}

\subsection{Locating a Research Gap}


\subsection{GNSS Spoofing}


% ====================================================================
\section{Evaluation Plan}

\subsection{Classroom Deployment}

\subsection{Expert Walkthrough}

\subsection{Task-Based Study}

\subsection{Case-Study Validation}

% ====================================================================
\section{Limitations}


% ====================================================================
\section{Conclusion}


% ====================================================================
\section*{Acknowledgment}
The authors thank \needdata{}.

\nocite{*}   % remove once \needcite{} markers become real \cite{} calls
\bibliographystyle{IEEEtran}
\bibliography{refs}

\end{document}
